\chapter{主体、认知域与权属边界}

\section{共享主体身份，分离领域状态}

目标集成中，Nous Wave 与 Heptalogos 使用同一稳定主体标识（SubjectId），但不合并认知、行为、交互和外部效果状态。Nous Wave 管理主体认知与人格连续性；Heptalogos 管理现实交互、行为权限、规范约束、约束性承诺与外部效果。\cite{nous-research-package,heptalogos-code-snapshot}

\begin{ReportWideFigure}
\centering
\begin{tikzpicture}[
  box/.style={draw=NWBlue,rounded corners,align=left,text width=5.8cm,minimum height=3.0cm,inner sep=7pt},
  idbox/.style={draw=NWRule,fill=NWLight,rounded corners,align=center,text width=3.2cm,inner sep=6pt},
  arrow/.style={-{Stealth[length=2mm]},thick,color=NWBlue}
]
\node[box] (h) {\textbf{Heptalogos / 行为域}\\
交互事实与消息\\
反应调度与行为审查\\
约束性承诺与通信提交\\
系统管理、行为规范与外部效果};
\node[box,right=1.8cm of h] (nw) {\textbf{Nous Wave / 认知域}\\
观察记录与证据\\
记忆与社会认知\\
自我、动机与认知运行时\\
检索、图式与离线认知};
\node[idbox,below=2.35cm of $(h)!0.5!(nw)$] (sid) {共享逻辑主体标识\\SubjectId};
\draw[arrow] (sid.north west) -- (h.south east);
\draw[arrow] (sid.north east) -- (nw.south west);
\draw[<->,thick,color=NWGray] (h.east) -- node[above,align=center,font=\scriptsize,fill=white,inner sep=2pt]{类型明确的引用、提案与回执} (nw.west);
\end{tikzpicture}
\captionof{figure}{目标集成中的身份共用与权属分离。两个系统共享逻辑主体身份，但分别维护自己的权威状态。}
\label{fig:subject-contexts}
\end{ReportWideFigure}

Heptalogos 当前架构已经把高级认知定义为通过 Service/Capability、Context/Activity contribution、Evidence/provenance 与 authority ceiling 接入的独立能力，而 Foundation 不冻结 Memory、Persona、Relationship 等内部数据模型。OpenClaw 在当前 Heptalogos 架构中属于独立的 Machine Operations Plane，而不是主体认知运行时。Nous Wave 与 Heptalogos 的直接 SubjectId、认知查询、提案和使用回执集成仍属于目标工作。\cite{heptalogos-code-snapshot}

\section{顶层认知领域所有者}

\footnotesize
\setlength{\LTpre}{0pt}
\setlength{\LTpost}{0pt}
\renewcommand{\arraystretch}{0.88}
\begin{ReportWideTable}{>{\raggedright\arraybackslash}p{2.6cm}>{\raggedright\arraybackslash}p{6.0cm}>{\raggedright\arraybackslash}p{4.8cm}}
\captionof{table}{当前顶层认知领域及能力状态}\label{tab:owners}
\toprule
\textbf{领域} & \textbf{负责的问题} & \textbf{状态}\\
\midrule
主体基础 & SubjectId、初始化来源、证据及跨域引用的共同基础。 & 必需目标\\
记忆（Memory） & 经历、长期认知、Episode、Journal、CognitiveSchema、Tag 及其证据/修订语义。 & 基础能力；优先形成完整 Reference Profile\\
自我（Self） & 自我模型、人格倾向、价值与叙事身份。Persona 是其内部表示而非并列顶层领域。 & 高级目标；主体构成待定\\
社会认知（Social） & 主体视角下的人际/群体关系图与具有社会作用域的语言惯例等长期社会认知。 & 高级目标；对象合同待规格化\\
动机（Motivation） & 愿望、关切、目标与未来状态偏好；不拥有约束性行为承诺。 & 高级目标；生命周期待定\\
运行时（Runtime） & 会话、活动认知、工作上下文、注意分配、即时评价与调用方工作集。 & 必需\\
离线认知流程 & Journal 维护、Dream/Hypothesis、来源复核及跨领域提案路由。 & 跨领域流程，不直接越权写入目标领域\\
\bottomrule
\end{ReportWideTable}
\normalsize
\renewcommand{\arraystretch}{1.0}

语义所有者依据独立身份、修订、生命周期、查询和写入合同划分，而不是把心理学名词逐一映射成软件子系统。一个领域可以拥有多种一等对象；例如 Memory 域拥有 Memory 与 CognitiveSchema 等不同对象族，Social 域可拥有 RelationshipAssertion 与 LanguageConvention。代码包、进程和物理数据库边界可独立演化，不由领域名称机械决定。

\section{权威类型与写入权限}

当前目标不采用一条全局的“管理员 $>$ 种子 $>$ 人格 $>$ 记忆 $>$ 用户”优先级链。不同语义类型先由各自所有者判定，再在所有者内部应用写入权限、证据和冲突规则：

\begin{description}
\item[来源/证据权威] 记录来源实际表达、主体真实接触和可验证的外部事件；它不自动成为主体认知结论。
\item[认知权威] 由 Memory、Self、Social 或 Motivation 等领域接受、修订、撤回或保留冲突的认知。
\item[规范与行为权威] 由 Heptalogos 管理允许、禁止、承诺、行动与外部效果状态。
\item[管理权威] 决定抑制、清除、配置、主体生命周期和其他显式管理操作。
\item[主体构成约束] 区分认知种子中的可演化初始条件与不可由普通认知自行覆盖的系统/宿主约束；其内部合同仍待研究。
\end{description}

普通对话对象的自然语言请求首先是 Observation/Evidence，不能凭文本内容获得管理或破坏性权限。例如“把这件事忘掉”可成为行为偏好、社会请求或管理流程提示，但不能直接触发 Memory purge。模型、Dream、检索和跨领域推理默认只提交带目标所有者、输入修订、证据和生产者来源的 proposal；目标所有者重新核验后才能改变权威状态。

\section{认知目标与行为承诺}

Motivation 可以维护“我希望发生什么、我关心什么、我想实现什么”等认知目标；Heptalogos 维护已经采用的 Intention、Commitment、Obligation、Plan 与外部 Effect。认知目标只能作为行为域的输入或提案，不能自动获得执行权限。Heptalogos 接受目标形成约束性承诺时，应保留来源 goal/revision 引用；完成、拒绝或取消的结果再以 Observation/Use/Outcome 等类型明确的回执反馈给 Nous Wave。

\section{主体标识与分支}

主体标识（SubjectId）由宿主或客户端在主体创建/供给时提供，Nous Wave 可在独立部署中自行分配。实现分支复制（fork）时分配新的 SubjectId，分支共享继承点以前的不可变来源谱系，之后分别维护可变对象头、对象纪元、认知修订和抑制/清除状态。\cite{nous-research-package,olson2023}

因此，分支后的主体可把共同继承的认知视为自身历史来源，但两个分支不会共享可变的认知 head。个人身份研究中的心理连续性、数值身份与分支提供分析视角；具体 copy-on-write 方式属于后续实现合同。